\documentclass[10pt,a4paper]{amsart}
\usepackage[margin=19mm]{geometry}
\usepackage{amsmath,amssymb,mathtools}
\usepackage[scr=rsfs,cal=euler]{mathalfa}
\usepackage{xfrac}
\usepackage[unicode,hidelinks,bookmarks=false]{hyperref}
\newcommand{\ie}{i.e.\ }
\newcommand{\cf}{cf.\ }
\newcommand{\resp}{respectively\ }
\newcommand{\Subsection}{\subsection}
\providecommand{\nobreakdash}{\nobreak\hbox{-}}
\setlength{\emergencystretch}{2em}
\multlinegap0pt
\usepackage{bm}
\usepackage{bbm}
\usepackage{dsfont}
\usepackage{xspace}
\usepackage{cancel}
\usepackage{mathtools}
\usepackage{amsthm}
\makeatletter
\@ifundefined{theorem}{\newtheorem{theorem}{Theorem}}{}
\makeatother
\title[The number of inscribed and circumscribed graphs of a convex]{The number of inscribed and circumscribed graphs of a convex polyhedron}
\author{Yagub N. Aliyev}
\date{Source: arXiv:2312.09552v8 (CC BY 4.0)}
\begin{document}
\maketitle

\noindent\emph{Source and licence.} The number of inscribed and circumscribed graphs of a convex polyhedron. Version arXiv:2312.09552v8 (CC BY 4.0), distributed under the Creative Commons Attribution 4.0 International licence, as recorded in the arXiv licence metadata for this exact version and shown on the abstract page https://arxiv.org/abs/2312.09552v8. Full rights notice: licenses/arxiv-2312.09552-CC-BY-4.0.txt.\par\smallskip

\noindent\textit{Editorial scope.} Counted unit: the Theorem whose source label is \texttt{None} in the article (source0.tex lines 228--230, source label \texttt{None}), delivered together with the article's own complete original proof of it (source0.tex lines 232--252, one \texttt{proof} environment of 21 lines). No mathematics is written, repaired or re-derived: statement and proof are the article's own text line for line. Required context retained uncounted: none. Every definition, display and label of those lines is retained unchanged. Omitted from this package and not redistributed: 1--227, 231--231, 253--588. Nothing inside a retained excerpt is omitted or altered: every retained body is delivered exactly as the article's lines are written, so no omission receipt of any kind accompanies this package. Citations: the retained excerpts carry no citation command at all, so no bibliography entry of the article is transcribed and no reference is redistributed. Reference adaptation: none was needed and none was made; every reference inside the delivered unit resolves inside it, so no unresolved reference or citation is produced. Printed slips kept literal: no printed word, symbol, hypothesis or number of the counted statement or of its proof was altered. Layout and class: the article's own class is \texttt{birkjour} with packages including hyperref, amsmath, amsthm, amssymb, enumitem, graphicx, caption, subcaption, fontenc, amsfonts, float, cancel; the wrapper is the lane's pinned \texttt{amsart} editorial wrapper at 10pt on A4, which restates the theorem-like environments the retained text needs and adds the article's own macro definitions that the retained text needs (none), with extra wrapper packages (bm, bbm, dsfont, xspace, cancel, mathtools). The delivered numbering is the unit's own. Assets: no retained span contains a figure environment, a graphics-inclusion command, a PDF-page inclusion, a table or a TikZ picture; no image file is copied into this package, the package declares no asset dependency and no image file is redistributed with it. Licence: version arXiv:2312.09552v8 of this article is distributed under the Creative Commons Attribution 4.0 International licence, as recorded in the arXiv licence metadata for this exact version and shown on the abstract page https://arxiv.org/abs/2312.09552v8. A full rights notice is delivered at licenses/arxiv-2312.09552-CC-BY-4.0.txt. Characters: every retained excerpt is pure ASCII, so the delivered engine, XeLaTeX with its default Latin Modern text font, reports no missing character.
\begin{theorem}
The number of convex polygons $B_0 B_1 B_2\ldots B_{n-1}$, that can be inscribed into a given convex polygon $A_0 A_1 A_2\ldots A_{n-1}$ and circumscribed around around another convex polygon $C_0 C_1 C_2\dots C_{n-1}$ is at most 4.
\end{theorem}

\begin{proof}
We saw in Example 1 (see also more general Example 2) that the case with 4 polygons is possible. Suppose on the contrary that there is a configuration with at least 5 polygons. Denote 5 of them by $B^j_0 B^j_1 B^j_2\ldots B^j_{n-1}$ $(j=1,2,\ldots,5)$ so that $|A_0B^j_0|<|A_0B^{j+1}_0|$, for $j=1,2,3,4$ (see Figure 7). Immediately, $|A_iB^j_i|<|A_iB^{j+1}_i|$, for all $i=0,1,\ldots,n-1$ and $j=1,2,3,4$. Consider the first two of the polygons: $B^1_0 B^1_1 B^1_2\ldots B^1_{n-1}$ and $B^{2}_0 B^{2}_1 B^{3}_2\ldots B^{3}_{n-1}$.
Then the vertices of $C_0 C_1 C_2\dots C_{n-1}$ are $C_i=B^1_iB^1_{i+1}\cap B^{2}_iB^{2}_{i+1}$ $(i=0,1,\ldots,n-1)$, where all the indices are taken modulo $n$. Note that $B^{3}_0 B^{3}_1$ can not pass through $C_0$. Indeed, otherwise $B^{3}_i B^{3}_{i+1}$ pass through $C_i$ for all $i=0,1,\ldots,n-1$, and by Lemma 2.1,
$$
\frac{|B^1_iB^2_i|}{|B^2_iB^3_i|}>\frac{|B^1_{i+1}B^2_{i+1}|}{|B^2_{i+1}B^3_{i+1}|},
$$
for all $i=0,1,\ldots,n-1$, which is impossible, because it implies $
\frac{|B^1_0B^2_0|}{|B^2_0B^3_0|}>\frac{|B^1_0B^2_0|}{|B^2_0B^3_0|}.
$
So, by convexity of $A_0 A_1 A_2\ldots A_{n-1}$ and $C_0 C_1 C_2\dots C_{n-1}$, $B^{3}_0 B^{3}_1$ passes through $C_1$. This means that $C_i\in B^{3}_{k+i} B^{3}_{k+i+1}$ $(i=0,1,\ldots,n-1)$ for $k=-1$. We can show that $k=0$ and $k=-1$ are the only possible choices of $k$ for which $C_i\in B^{j}_{k+i} B^{j}_{k+i+1}$ ($i=0,1,\ldots,n-1$, and $j=1,2,3,4$).
Indeed, if $n>3$, then all sets $$\mathbb{M}_i=\triangle A_iA_{i+1}A_{i+2}\cap \triangle A_{i-1}A_iA_{i+1}\ (i=0,1,\ldots,n-1)$$ are disjoint (for $\mathbb{M}_i$ only the interior regions of the triangles are considered) and $C_i\in \mathbb{M}_i$. If $n=3$, then the previous argument fails because $\triangle A_iA_{i+1}A_{i+2}$ $(i=0,1,2)$ coincides with $\triangle A_0A_{1}A_{2}$ and therefore $\mathbb{M}_i=\triangle A_0A_{1}A_{2}$. 
But if $n=3$, then the only remaining option for $k$ is $k=1$ modulo 3, and this case is impossible because in this case there is $j\ne 1,2$ such that $C_0\in B^{j}_{1} B^{j}_{2}$, and therefore $B^{j}_{1}$ is closer to $A_1$ than $B^{1}_{1}$, which is against our initial assumptions.

Consequently, $B^{4}_{0} B^{4}_1$, $B^{5}_{0} B^{5}_1$ also pass through $C_{1}$, and therefore $B^{3}_i B^{3}_{i+1}$, $B^{4}_i B^{4}_{i+1}$, and $B^{5}_i B^{5}_{i+1}$ pass through $C_{i+1}$ for all $i=0,1,\ldots,n-1$.
By Lemma 2.2,
$$
\frac{|B^3_{i}B^4_{i}|}{|B^4_{i}B^5_{i}|}>\frac{|B^3_{i+1}B^4_{i+1}|}{|B^4_{i+1}B^5_{i+1}|},
$$
for all $i=0,1,\ldots,n-1$, which is again impossible.
This shows that our assumption about the existence of 5 polygons was wrong and therefore the maximal number of polygons is 4.
\end{proof}

\end{document}
